\section{Discussion}
By analysing the final quantitative and qualitative evaluation results presented in Table \ref{tab:probe_comparison}, several insights can be drawn regarding the platform's suitability for a university-level computer systems course.

\subsection{Evaluation against Initial Pedagogical and Technical Goals}

Our first goal was affordability, continuing the cost-conscious design philosophy inherited from the platform's initial development iteration \cite{kelhala:2024}. Demonstrating exceptional cost-efficiency, the second core software approach (\texttt{pico-debug}) is the most cost-effective solution, as it is completely free and requires no wiring or external hardware \cite{pico_debug}. For physical hardware debugging, our custom Second Pico approach remains highly affordable, costing only 16€ (including the 3D-printed casing and wires), which is well within our affordability goal. Both methods are inexpensive compared to the industry-standard SEGGER J-Link, which sits at 77--89€ per unit at scale and 800€ individually for large-scale classroom deployment \cite{jlink}.

Our second goal was ease of setup and reproducibility. By utilising the VM and Docker container environments, we minimised the common host-side configuration errors discussed by Kelhälä et al. \cite{kelhala:2024}. Specifically, we resolved the limitation of the VM's heavy system overhead on older student laptops by implementing the Docker container, which achieves faster startup times whilst utilising process-level isolation rather than full hardware emulation \cite{rosen_namespaces_slides}. Nevertheless, we still maintain the Virtual Machine option as it has historically been the established standard for the Computer Systems course and remains a highly usable, stable solution \cite{tkj_vm_history}.

Our third goal was full debugging functionality, particularly multi-core Symmetric Multiprocessing (SMP) and FreeRTOS task awareness. Whilst the virtual \texttt{pico-debug} software probe represents a highly elegant, zero-cost alternative, it permanently occupies the target's second CPU core \cite{pico_debug}, preventing students from writing or debugging dual-core concurrent programs and conflicting with standard USB serial outputs (such as the Pico SDK's \texttt{stdio\_init\_all()}). Consequently, the Second Pico hardware probe is the only affordable solution that preserves the microcontroller's full dual-core architecture whilst supporting native FreeRTOS thread-awareness within the debugging pipeline.

\subsection{Relation to Related Work}

In comparison to standard, un-enclosed hardware configurations described in general educational literature, such as those by Harvey \cite{harvey_pico}, or the original, bare TKJHAT printed circuit board (PCB) \cite{kelhala:2024}, our platform introduces a physical, custom-designed 3D-printed casing. Whilst exposed, un-enclosed circuit boards are highly susceptible to electrostatic discharge (ESD), physical damage, and accidental short circuits on laboratory desks, our modular casing integrates the Pico and the custom TKJHAT board into a secure, robust 3D-printed unit.

Furthermore, our platform directly addresses three critical limitations of previous iterations of the Oulu Computer Systems course. Firstly, we successfully resolved the long-standing lack of a proper, standardised debugging solution, transitioning students away from primitive print-based troubleshooting \cite{tkj_vm_history}. Secondly, to support resource-constrained students whose laptops struggle with the heavy CPU and memory overhead of traditional virtual machines, we devised a lightweight, high-performance Docker container accessed via browser-based \texttt{code-server} \cite{code_server}. Finally, we resolved the problem of having multiple, uncoordinated user interfaces spread across different VS Code extensions by developing our custom, unified \texttt{Pico-Debug-Dashboard} to provide a clean, single-click build and debug interface \cite{pico_debug_dashboard}.

\subsection{Educational Recommendations for Lab Deployment}
To optimise the educational utility, student onboarding, and hardware resource management, we recommend a tiered deployment strategy aligned directly with the Lovelace course structure:
\begin{itemize}
    \item \textbf{Primary Development Environment:} The standardised Virtual Machine (via VirtualBox OVA) and Docker container environments, fully integrated with our custom \texttt{Pico-Debug-Dashboard} Visual Studio Code extension, should be deployed as the official development environments for both courses \cite{pico_debug_dashboard}. This ensures that every student operates within a consistent, pre-configured software environment, completely eliminating host-side setup, hardware-induced restrictions, compilation errors across all modules \cite{kelhala:2024}.
    
    \item \textbf{Mass Deployment for Lovelace Exercises (pico-debug):} For standard, individual Lovelace exercises (Modules 1--3) and initial mass laboratory exercises, the virtual \texttt{pico-debug} (Core) software setup should be utilised as the primary debugger \cite{pico_debug}. Because it requires no hardware assembly, zero wiring, and zero cost, it is highly scalable for mass deployment, allowing all students to perform basic step-through execution and register inspection with zero hardware friction.
    
    \item \textbf{Advanced Troubleshooting and Course Projects (Second Pico):} When students encounter highly complex bugs, require advanced execution tracing, or begin developing their collaborative multi-threaded FreeRTOS group projects in Module 2 and Module 4, they should transition to the Second Pico setup \cite{pico_getting_started}. This physical debugger is necessary to preserve the microcontroller's second CPU core, enable native FreeRTOS task-awareness, and provide an isolated hardware reset mechanism if the target's program crashes or hangs during laboratory sessions \cite{pico_getting_started}.
    \item \textbf{Standardised Hardware Enclosure (3D-Printed Casing):} To protect the physical components in a busy university laboratory, we recommend enclosing the target hardware inside our custom, modular 3D-printed casing. The enclosure is designed to physically accommodate and completely cover every current iteration of the custom educational HAT used in the course, providing uniform, robust protection against accidental short-circuits and mechanical damage. 
\end{itemize}
\section{Conclusion}
Based on the quantitative and qualitative data presented in Table \ref{tab:probe_comparison}, the following conclusions can be drawn:

The \textbf{SEGGER J-Link}, while representing the industry standard for professional embedded development, was found to be the least suitable for the course environment due to its prohibitive cost (77--89 €) at scale and by the further complexity of its proprietary software stack. Conversely, the \textbf{pico-debug (Core)} implementation represents the most efficient solution in terms of cost and physical footprint. However, its low maintenance score and the archived status of its RP2040 firmware introduce a significant risk for a curriculum intended to remain operational for several years.

The \textbf{Second Pico} implementation achieved the highest total qualitative score (12.3/15). This methodology offers a "middle ground" that provides the high-level maintenance support of an official Raspberry Pi product while remaining affordable for large-scale classroom deployment. Its ability to provide true dual-core debugging (SMP support) without the maintenance risks associated with virtual probes makes it the most robust choice for a multi-year educational programme.

Consequently, while \textit{pico-debug} serves as an excellent zero-cost entry point for individual students, the Second Pico implementation is recommended as the primary debugging methodology for the Computer Systems course. When integrated with the custom IDE dashboard and 3D-printed casing, it provides the most balanced, reliable, and "snappy" user experience for incoming students.

The primary \textbf{tangible outputs} developed and delivered as part of this project are fully open-source and include:
\begin{itemize}
    \item \textbf{The Virtual Machine Environment:} A pre-configured, standardised Debian 13 (Trixie) virtual appliance containing the complete compiler toolchain, SDKs, and Visual Studio Code: \\ \url{https://drive.google.com/file/d/1FaViwvn3uLcmR33f-7BOAskOCEKr74pN/view?usp=drive_link} \cite{tkj_vm_history}.
    \item \textbf{The Docker Environment \& IDE Dashboard:} The official GitHub repository containing our custom-built, browser-based \texttt{code-server} container and the TypeScript-based \texttt{Pico-Debug-Dashboard} extension: \\ \url{https://github.com/CenderVN/Pico-Debug-Dashboard} \cite{pico_debug_dashboard}.
    \item \textbf{The 3D-Printed Case CAD Models:} The official GitHub repository containing the 3D-printable `.step` and `.stl` files for our modular, protective case: \\ \url{https://github.com/PrabowoXYZ/TKJHAT_Case}.
\end{itemize}

Ultimately, we hope the implementation of this platform will have a significant positive impact on embedded systems education at the University of Oulu. By enabling students with lower-end laptops to have a much smoother development experience, protecting laboratory hardware, and simplifying the debugging interface, this platform allows future students, teaching assistants, and course instructors to focus entirely on the core computer systems concepts, making embedded systems education more accessible, engaging, and rewarding.